% GENERATED FILE. Edit canonical lessons, then run npm run generate:books.
% canonical-source-hash: fnv1a64:49cc159c80f4d711
% canonical-lessons: LA-C310-iuro, LA-C310-periurium, LA-C310-mentior, LA-C310-arguo, LA-C310-indicium

\chapter{Vērum aut falsum}
\label{ch:la-verum-aut-falsum}
\hlchaptermodality{310}

\begin{chapteropening}
\textbf{By the end of this chapter:} \emph{I can talk about truth and lies in court: to swear, perjury, to lie, to prove and information.}

Complete the last lesson of chapter 310: I can talk about truth and lies in court: to swear, perjury, to lie, to prove and information.
\end{chapteropening}

\section[iūrō, iūrāre]{iūrō, iūrāre — to swear (an oath)}

\label{lesson:LA-C310-iuro}



\begin{quote}
Before the new one: say the Latin for a patron, then the Latin for evidence.
\end{quote}

\subsection*{You'll want to know: iūrō, iūrāre}
\textbf{iūrō, iūrāre} — \textquotedblleft{}to swear (an oath)\textquotedblright{}.

\begin{quote}
Testis iūrāvit sē vēra dīcere.
\end{quote}

\begin{quote}
\emph{The witness swore that he was telling the truth.}
\end{quote}

\emph{iūrō, iūrāre, iūrāvī, iūrātum}, with accusative and infinitive. English \emph{jury} comes from it through French.

\begin{cousinweb}
\textbf{iūrō} is built on \emph{iūs}, \textquotedblleft{}right, law\textquotedblright{} (EtymDB): to bind oneself by law. English \emph{jury} comes from it through French.
\end{cousinweb}

\begin{grammarlens}[title={What you've built}]
One more word for telling, explaining and describing at B1.
\end{grammarlens}

\subsection*{Your turn}
\begin{itemize}
\raggedright
  \item \emph{Say it:} \emph{iūrō, iūrāre}
  \item \emph{Say it:} \emph{iūrō, iūrāre}, once more
  \item \emph{Say it:} a sentence of your own with \emph{iūrō, iūrāre}
  \item \emph{Recall:} say the Latin for a patron, then the Latin for evidence, then say \emph{iūrō, iūrāre} again
\end{itemize}

\begin{tcolorbox}[breakable,colback=teal!4,colframe=teal!35!black,title={Before you move on}]
What does iūrō, iūrāre mean? (\textquotedblleft{}To swear (an oath)\textquotedblright{}.) Say it once more.
\end{tcolorbox}
\section[periūrium, periūriī]{periūrium, periūriī — perjury, a false oath}

\label{lesson:LA-C310-periurium}



\begin{quote}
Before the new one: say the Latin for evidence, then the Latin for to swear.
\end{quote}

\subsection*{You'll want to know: periūrium, periūriī}
\textbf{periūrium, periūriī} — \textquotedblleft{}perjury, a false oath\textquotedblright{}.

\begin{quote}
Periūrium grave crīmen erat.
\end{quote}

\begin{quote}
\emph{Perjury was a serious crime.}
\end{quote}

Neuter, from \emph{periūrus}, 'oath-breaking': \emph{per-} here means 'wrongly'. English \emph{perjury}.

\begin{cousinweb}
\textbf{periūrium} is built on \emph{periūrus}, \textquotedblleft{}oath-breaking\textquotedblright{} (EtymDB), from \emph{per-}, here \textquotedblleft{}wrongly\textquotedblright{}, and the root of \emph{iūrō}, \textquotedblleft{}to swear\textquotedblright{}. English \emph{perjury} comes from it.
\end{cousinweb}

\begin{grammarlens}[title={What you've built}]
One more word for telling, explaining and describing at B1.
\end{grammarlens}

\subsection*{Your turn}
\begin{itemize}
\raggedright
  \item \emph{Say it:} \emph{periūrium, periūriī}
  \item \emph{Say it:} \emph{periūrium, periūriī}, once more
  \item \emph{Say it:} a sentence of your own with \emph{periūrium, periūriī}
  \item \emph{Recall:} say the Latin for evidence, then the Latin for to swear, then say \emph{periūrium, periūriī} again
\end{itemize}

\begin{tcolorbox}[breakable,colback=teal!4,colframe=teal!35!black,title={Before you move on}]
What does periūrium, periūriī mean? (\textquotedblleft{}Perjury, a false oath\textquotedblright{}.) Say it once more.
\end{tcolorbox}
\section[mentior, mentīrī]{mentior, mentīrī — to lie, to tell a lie}

\label{lesson:LA-C310-mentior}



\begin{quote}
Before the new one: say the Latin for to swear, then the Latin for perjury.
\end{quote}

\subsection*{You'll want to know: mentior, mentīrī}
\textbf{mentior, mentīrī} — \textquotedblleft{}to lie, to tell a lie\textquotedblright{}.

\begin{quote}
Servus mentītus est; nam pōculum ipse frēgerat.
\end{quote}

\begin{quote}
\emph{The slave lied, for he had broken the cup himself.}
\end{quote}

\emph{mentior, mentīrī, mentītus sum}, a deponent. Its noun \emph{mendācium}, 'a lie', is taught.

\begin{cousinweb}
\textbf{mentior} is probably built, says Lewis \& Short, on the root of \emph{mēns}, \textquotedblleft{}mind\textquotedblright{}, with a first sense \textquotedblleft{}to invent, to make up\textquotedblright{}; this is not certain.
\end{cousinweb}

\begin{grammarlens}[title={What you've built}]
One more word for telling, explaining and describing at B1.
\end{grammarlens}

\subsection*{Your turn}
\begin{itemize}
\raggedright
  \item \emph{Say it:} \emph{mentior, mentīrī}
  \item \emph{Say it:} \emph{mentior, mentīrī}, once more
  \item \emph{Say it:} a sentence of your own with \emph{mentior, mentīrī}
  \item \emph{Recall:} say the Latin for to swear, then the Latin for perjury, then say \emph{mentior, mentīrī} again
\end{itemize}

\begin{tcolorbox}[breakable,colback=teal!4,colframe=teal!35!black,title={Before you move on}]
What does mentior, mentīrī mean? (\textquotedblleft{}To lie, to tell a lie\textquotedblright{}.) Say it once more.
\end{tcolorbox}
\section[arguō, arguere]{arguō, arguere — to prove; to accuse}

\label{lesson:LA-C310-arguo}



\begin{quote}
Before the new one: say the Latin for perjury, then the Latin for to lie.
\end{quote}

\subsection*{You'll want to know: arguō, arguere}
\textbf{arguō, arguere} — \textquotedblleft{}to prove; to accuse\textquotedblright{}.

\begin{quote}
Haec epistula eum mentītum esse arguit.
\end{quote}

\begin{quote}
\emph{This letter proves that he lied.}
\end{quote}

\emph{arguō, arguere, arguī, argūtum}: to make clear, and so to prove or to charge. English \emph{argue}.

\begin{cousinweb}
\textbf{arguō} first meant \textquotedblleft{}to make clear\textquotedblright{}; Lewis \& Short connect it with the root of \emph{argentum}, \textquotedblleft{}silver\textquotedblright{}, and Greek \emph{argos}, \textquotedblleft{}bright\textquotedblright{}. English \emph{argue} comes from it through French.
\end{cousinweb}

\begin{grammarlens}[title={What you've built}]
One more word for telling, explaining and describing at B1.
\end{grammarlens}

\subsection*{Your turn}
\begin{itemize}
\raggedright
  \item \emph{Say it:} \emph{arguō, arguere}
  \item \emph{Say it:} \emph{arguō, arguere}, once more
  \item \emph{Say it:} a sentence of your own with \emph{arguō, arguere}
  \item \emph{Recall:} say the Latin for perjury, then the Latin for to lie, then say \emph{arguō, arguere} again
\end{itemize}

\begin{tcolorbox}[breakable,colback=teal!4,colframe=teal!35!black,title={Before you move on}]
What does arguō, arguere mean? (\textquotedblleft{}To prove; to accuse\textquotedblright{}.) Say it once more.
\end{tcolorbox}
\section[indicium, indiciī]{indicium, indiciī — information, evidence; a sign}

\label{lesson:LA-C310-indicium}



\begin{quote}
Before the new one: say the Latin for to lie, then the Latin for to prove.
\end{quote}

\subsection*{You'll want to know: indicium, indiciī}
\textbf{indicium, indiciī} — \textquotedblleft{}information, evidence; a sign\textquotedblright{}.

\begin{quote}
Servus cōnsulī indicium dedit.
\end{quote}

\begin{quote}
\emph{A slave gave information to the consul.}
\end{quote}

Neuter. \emph{Indicium facere}: to give information. A clue in a story is also an \emph{indicium}.

\begin{cousinweb}
\textbf{indicium} is built on \emph{index}, \textquotedblleft{}one who points out, an informer\textquotedblright{} (EtymDB), and so on the verb meaning \textquotedblleft{}to point out\textquotedblright{}: a pointing out.
\end{cousinweb}

\begin{grammarlens}[title={What you've built}]
One more word for telling, explaining and describing at B1.
\end{grammarlens}

\subsection*{Your turn}
\begin{itemize}
\raggedright
  \item \emph{Say it:} \emph{indicium, indiciī}
  \item \emph{Say it:} \emph{indicium, indiciī}, once more
  \item \emph{Say it:} a sentence of your own with \emph{indicium, indiciī}
  \item \emph{Recall:} say the Latin for to lie, then the Latin for to prove, then say \emph{indicium, indiciī} again
\end{itemize}

\begin{tcolorbox}[breakable,colback=teal!4,colframe=teal!35!black,title={Before you move on}]
What does indicium, indiciī mean? (\textquotedblleft{}Information, evidence; a sign\textquotedblright{}.) Say it once more.
\end{tcolorbox}
